\documentclass{amsart}
% For dealing with references we use the comment environment
\usepackage{verbatim}
\newenvironment{reference}{\comment}{\endcomment}
%\newenvironment{reference}{}{}
\newenvironment{slogan}{\comment}{\endcomment}
\newenvironment{history}{\comment}{\endcomment}

% For commutative diagrams we use Xy-pic
\usepackage[all]{xy}

% We use 2cell for 2-commutative diagrams.
\xyoption{2cell}
\UseAllTwocells

% We use multicol for the list of chapters between chapters
\usepackage{multicol}

% This is generally recommended for better output
\usepackage{lmodern}
\usepackage[T1]{fontenc}

% For cross-file-references

% Package for hypertext links:
\usepackage{hyperref}

% For any local file, say "hello.tex" you want to link to please

% Theorem environments.
%
\theoremstyle{plain}
\newtheorem{theorem}[subsection]{Theorem}
\newtheorem{proposition}[subsection]{Proposition}
\newtheorem{lemma}[subsection]{Lemma}

\theoremstyle{definition}
\newtheorem{definition}[subsection]{Definition}
\newtheorem{example}[subsection]{Example}
\newtheorem{exercise}[subsection]{Exercise}
\newtheorem{situation}[subsection]{Situation}

\theoremstyle{remark}
\newtheorem{remark}[subsection]{Remark}
\newtheorem{remarks}[subsection]{Remarks}

\numberwithin{equation}{subsection}

% Macros
%
\def\lim{\mathop{\mathrm{lim}}\nolimits}
\def\colim{\mathop{\mathrm{colim}}\nolimits}
\def\Spec{\mathop{\mathrm{Spec}}}
\def\Hom{\mathop{\mathrm{Hom}}\nolimits}
\def\Ext{\mathop{\mathrm{Ext}}\nolimits}
\def\SheafHom{\mathop{\mathcal{H}\!\mathit{om}}\nolimits}
\def\SheafExt{\mathop{\mathcal{E}\!\mathit{xt}}\nolimits}
\def\Sch{\mathit{Sch}}
\def\Mor{\mathop{\mathrm{Mor}}\nolimits}
\def\Ob{\mathop{\mathrm{Ob}}\nolimits}
\def\Sh{\mathop{\mathit{Sh}}\nolimits}
\def\NL{\mathop{N\!L}\nolimits}
\def\CH{\mathop{\mathrm{CH}}\nolimits}
\def\proetale{{pro\text{-}\acute{e}tale}}
\def\etale{{\acute{e}tale}}
\def\QCoh{\mathit{QCoh}}
\def\Ker{\mathop{\mathrm{Ker}}}
\def\Im{\mathop{\mathrm{Im}}}
\def\Coker{\mathop{\mathrm{Coker}}}
\def\Coim{\mathop{\mathrm{Coim}}}

% Boxtimes
%
\DeclareMathSymbol{\boxtimes}{\mathbin}{AMSa}{"02}

%
% Macros for moduli stacks/spaces
%
\def\QCohstack{\mathcal{QC}\!\mathit{oh}}
\def\Cohstack{\mathcal{C}\!\mathit{oh}}
\def\Spacesstack{\mathcal{S}\!\mathit{paces}}
\def\Quotfunctor{\mathrm{Quot}}
\def\Hilbfunctor{\mathrm{Hilb}}
\def\Curvesstack{\mathcal{C}\!\mathit{urves}}
\def\Polarizedstack{\mathcal{P}\!\mathit{olarized}}
\def\Complexesstack{\mathcal{C}\!\mathit{omplexes}}
% \Pic is the operator that assigns to X its picard group, usage \Pic(X)
% \Picardstack_{X/B} denotes the Picard stack of X over B
% \Picardfunctor_{X/B} denotes the Picard functor of X over B
\def\Pic{\mathop{\mathrm{Pic}}\nolimits}
\def\Picardstack{\mathcal{P}\!\mathit{ic}}
\def\Picardfunctor{\mathrm{Pic}}
\def\Deformationcategory{\mathcal{D}\!\mathit{ef}}

\setlength{\emergencystretch}{3em}
\begin{document}
\title[Stacks Project, Tag 01YQ]{368 --- Affineness descends under finite surjective morphisms to Noetherian schemes}
\maketitle
\noindent Copyright (C) 2005--2025 Johan de Jong. Stacks Project authors. Source: \href{https://stacks.math.columbia.edu/tag/01YQ}{Tag 01YQ}.

\noindent Editorial difficulty note (not author text). Difficulty H2: The finite cover need not be flat, so ordinary faithfully flat descent of affineness does not apply. The proof constructs coherent test sheaves above integral support strata and verifies cohomological devissage for all their ideal modifications, forcing H1 vanishing on the target..

\noindent Editorial provenance note (not author text). History: excerpt from revision \texttt{a0444\allowbreak 6e57e\allowbreak c1fbc\allowbreak 252a8\allowbreak 71afc\allowbreak ec775\allowbreak 2fb28\allowbreak 07b14}, coherent.tex, lines 3556--3617. Reference presentation was adapted; original-author context and core support blocks were retained or added. The mathematical wording is unchanged. Licensed under GNU FDL 1.2 or later; no invariant sections or cover texts. See ../licenses/COPYING. Context blocks reproduce author definitions and supporting results; they are not additional counted problems.

\begin{lemma}
\label{lemma-image-affine-finite-morphism-affine-Noetherian}
Let $f : Y \to X$ be a morphism of schemes.
Assume
\begin{enumerate}
\item $f$ finite,
\item $f$ surjective,
\item $Y$ affine, and
\item $X$ Noetherian.
\end{enumerate}
Then $X$ is affine.
\end{lemma}

\begin{proof}
We will prove that under the assumptions of the lemma for any coherent
$\mathcal{O}_X$-module $\mathcal{F}$ we have $H^1(X, \mathcal{F}) = 0$.
This will in particular imply that $H^1(X, \mathcal{I}) = 0$
for every quasi-coherent sheaf of ideals of $\mathcal{O}_X$. Then it
follows that $X$ is affine from either
Lemma \href{https://stacks.math.columbia.edu/tag/01XF}{lemma quasi compact h1 zero covering (Tag 01XF)} or
Lemma \href{https://stacks.math.columbia.edu/tag/01XG}{lemma quasi separated h1 zero covering (Tag 01XG)}.

\medskip\noindent
Let $\mathcal{P}$ be the property of coherent sheaves
$\mathcal{F}$ on $X$ defined by the rule
$$
\mathcal{P}(\mathcal{F}) \Leftrightarrow H^1(X, \mathcal{F}) = 0.
$$
We are going to apply Lemma \href{https://stacks.math.columbia.edu/tag/01YM}{lemma property higher rank cohomological (Tag 01YM)}.
Thus we have to verify (1), (2) and (3) of that lemma for $\mathcal{P}$.
Property (1) follows from the long exact cohomology sequence associated
to a short exact sequence of sheaves. Property (2) follows since
$H^1(X, -)$ is an additive functor. To see (3) let $Z \subset X$ be
an integral closed subscheme with generic point $\xi$.
Let $\mathcal{F}$ be a coherent sheaf on $Y$ such that
the support of $f_*\mathcal{F}$ is equal to $Z$
and $(f_*\mathcal{F})_\xi$ is annihilated by $\mathfrak m_\xi$,
see Lemma \href{https://stacks.math.columbia.edu/tag/01YO}{lemma finite morphism Noetherian (Tag 01YO)}. We claim that
taking $\mathcal{G} = f_*\mathcal{F}$ works. We only have to verify
part (3)(c) of Lemma \href{https://stacks.math.columbia.edu/tag/01YM}{lemma property higher rank cohomological (Tag 01YM)}.
Hence assume that $\mathcal{J} \subset \mathcal{O}_X$ is a
quasi-coherent sheaf of ideals such that
$\mathcal{J}_\xi = \mathcal{O}_{X, \xi}$.
A finite morphism is affine hence by
Lemma \href{https://stacks.math.columbia.edu/tag/01YP}{lemma affine morphism projection ideal (Tag 01YP)} we see that
$\mathcal{J}\mathcal{G} = f_*(f^{-1}\mathcal{J}\mathcal{F})$.
Also, as pointed out in the proof of
Lemma \href{https://stacks.math.columbia.edu/tag/01YP}{lemma affine morphism projection ideal (Tag 01YP)} the sheaf
$f^{-1}\mathcal{J}\mathcal{F}$ is a quasi-coherent $\mathcal{O}_Y$-module.
Since $Y$ is affine we see that $H^1(Y, f^{-1}\mathcal{J}\mathcal{F}) = 0$,
see Lemma \href{https://stacks.math.columbia.edu/tag/01XB}{lemma quasi coherent affine cohomology zero (Tag 01XB)}.
Since $f$ is finite, hence affine, we see that
$$
H^1(X, \mathcal{J}\mathcal{G}) =
H^1(X, f_*(f^{-1}\mathcal{J}\mathcal{F})) =
H^1(Y, f^{-1}\mathcal{J}\mathcal{F}) = 0
$$
by Lemma \href{https://stacks.math.columbia.edu/tag/089W}{lemma relative affine cohomology (Tag 089W)}.
Hence the quasi-coherent subsheaf $\mathcal{G}' = \mathcal{J}\mathcal{G}$
satisfies $\mathcal{P}$. This verifies property (3)(c) of
Lemma \href{https://stacks.math.columbia.edu/tag/01YM}{lemma property higher rank cohomological (Tag 01YM)} as desired.
\end{proof}
\end{document}
